\section{Methodology}
\label{sec:methods}

\subsection{Study Design}

The validation proceeds in three tiers of increasing empirical complexity. The first tier is a \textbf{theoretical/numerical foundation}: an idealised probit simulation (\cref{sec:results}; \texttt{run\_pef\_idealised\_probit\_sim.m}) that generates bivariate-normal $(X_{\mathrm{A}},X_{\mathrm{B}})$ on a $(\kappa,\rho,\delta/\sigma_{\mathrm{A}})$ grid under assumptions (A1)--(A2) and confirms the closed-form mapping from PEF to information content (\cref{eq:mi_closed}) in a controlled setting where those assumptions hold by construction (full model specification and validation in Supplementary~Note~S2). Because the sports KPIs analysed here are well approximated by bivariate normality (normality testing, \cref{sec:qc}), this idealisation provides an informative reference for the empirical data rather than a disconnected model. The second tier is the \textbf{primary empirical analysis} on professional sports KPIs. Here the full KPI inventory ($\PEFtotalStudies{}$ studies: $\PEFrugbyNkpis{}$ rugby, $\PEFfootNkpis{}$ football) provides a \emph{landscape characterisation} of where real metrics sit on the $(\kappa,\rho)$ plane (\cref{fig:pef_landscape,fig:pef_ml}; Supplementary Figures~S2--S3), whilst \emph{mechanistic confirmation} rests on four single-KPI exemplars---one per quadrant---with comparable signal strength (\cref{sec:exemplars,tab:exemplars}). The third tier is \textbf{supporting validation} on published data from four additional domains, assessing generalisability. Across these tiers the objectives were to: (1)~validate the PEF--information content relationship (\cref{eq:mi_closed}), (2)~confirm quadrant-specific predictions and the role of signal strength $\delta/\sigma_{\mathrm{A}}$ as a moderator, (3)~characterise the efficiency--power tension in Quadrant~4, and (4)~test cross-domain applicability.

\subsection{Primary Data: Sports KPIs}

The primary analysis uses match-level KPI data from two professional leagues:
\begin{itemize}
  \item \textbf{Rugby union.} United Rugby Championship ({URC}): $n=\PEFrugbyNgames$ games pooled across seasons 23/24 and 24/25 ($m=16$ teams, double round-robin). KPIs include rucks won, kick metres, kicks from hand, carries, penalties conceded, turnovers won, and others (\cref{fig:pef_landscape}).
  \item \textbf{Football.} English Championship: $n=\PEFfootNgames$ games pooled across the same two seasons ($m=24$ teams, double round-robin). KPIs include duels, cards, pressures, interceptions, fouls, and shots (\cref{fig:pef_landscape}). Association-football score and outcome modelling has a long econometric lineage \parencite{dixon1997,boulier2003}, with contemporary machine-learning approaches incorporating domain knowledge \parencite{berrar2019}.
\end{itemize}
Head-to-head matchups provide natural pairing through shared environmental conditions (weather, venue, officiating), and KPI variability across matches is well documented in professional sport \parencite{malcata2014}. Those rugby union lines of work motivate treating head-to-head KPIs as a primary testbed for the present framework \parencite{bennett2019descriptive,bennett2021predicting,scott2023urc,scott2023womens}. The sports data are well suited to testing the framework because the KPIs span a range of correlation values (including negative correlation from competitive dynamics) and variance ratios, populating all four quadrants.  To assess temporal stability, we compute $(\hat{\kappa},\hat{\rho})$ for each KPI in both seasons separately; \cref{fig:pef_landscape,fig:info_surface} show the four confirmatory exemplars of \cref{tab:exemplars} with season~23/24--24/25 drift segments, whilst Supplementary~\cref{fig:si_kpi_labelled} visualises the full KPI cloud.

\subsection{Supporting Data: Cross-Domain Validation}

To assess generalisability beyond sports, we applied the framework to four additional domains. Domain-level summaries are reported in \cref{tab:validation}; mean $(\kappa,\rho)$ positions for the supporting tier appear as overlay markers in \cref{fig:pef_landscape}.

\textbf{Healthcare.} Paired systolic ({BPXOSY1}) and diastolic ({BPXODI1}) blood pressure readings from the same oscillometric measurement event, drawn from the {NHANES} 2017 public-use examination file \texttt{P\_BPXO} \parencite{nhanes2017pbxo}. Each participant contributes one (systolic, diastolic) pair; pulse pressure $\mathrm{PP}=\mathrm{SBP}-\mathrm{DBP}$ is the relative feature, an established cardiovascular risk biomarker \parencite{franklin1999,blacher1999}. Observations with $\mathrm{SBP}\le\mathrm{DBP}$ were excluded as physiologically implausible ($N=10{,}352$ retained).

\textbf{Clinical Genomics.} Normalised {RNA}-seq gene-level expression from {GEO} series {GSE47462}: patient-matched normal breast tissue vs early neoplasia \parencite{geoGSE47462,brunner2014earlybreast}. Pairs were defined by shared subject identifiers in the series sample annotations; per-gene PEF summaries aggregate correlation and variance-ratio structure across those paired specimens.

\textbf{Finance.} Market data from Yahoo Finance across multiple asset classes and periods. Market-adjusted and benchmark-relative returns provide natural pairing \parencite{sharpe1966,sharpe1994,fama1970,carhart1997,fama1993}.

\textbf{Manufacturing.} Real industrial process telemetry from the {Bosch} production-line {CNC} dataset on {OpenML} (dataset id 752) \parencite{openml752bosch}. For each production cycle (one row), paired vectors are formed axis-wise: angular position $\theta_i$ versus angular velocity $\dot\theta_i$ ($i=1,\ldots,6$) and torque $\tau_j$ versus displacement metric $\mathrm{dm}_j$ ($j=1,\ldots,5$), yielding two PEF estimates per cycle. Summaries aggregate $N=16{,}384$ such pairings (8{,}192 cycles $\times$ two pairing types). The {NASA} {C-MAPSS} turbofan simulation \parencite{saxena2008} was used in earlier drafts but is not the default supporting manufacturing source here because it is simulated rather than observational plant data.

\subsection{Parameter Estimation}

For each dataset, the variance ratio and correlation were estimated using standard sample estimators
\begin{equation}
  \hat{\kappa} = \frac{s^{2}_{\mathrm{B}}}{s^{2}_{\mathrm{A}}},
  \quad
  \hat{\rho} =
  \frac{\sum_{i=1}^{n}(X_{\mathrm{A},i}-\bar{X}_{\mathrm{A}})(X_{\mathrm{B},i}-\bar{X}_{\mathrm{B}})}
       {\sqrt{\sum_{i=1}^{n}(X_{\mathrm{A},i}-\bar{X}_{\mathrm{A}})^{2}
              \sum_{i=1}^{n}(X_{\mathrm{B},i}-\bar{X}_{\mathrm{B}})^{2}}}.
\end{equation}
Under bivariate normality (Assumption~(A1)), these are maximum likelihood estimators. PEF was computed by substituting into \cref{eq:pef}.

\subsection{Empirical Units and Outcome Definitions}
\label{sec:outcome_defs}

Each sports KPI constitutes one \emph{study}: match-level home and away measurements pooled across seasons 23/24 and 24/25, yielding a single $(\hat\kappa,\hat\rho,\hat\eta)$ estimate per KPI. The full inventory comprises $\PEFtotalStudies{}$ studies ($\PEFrugbyNkpis{}$ rugby union, $\PEFfootNkpis{}$ football). Per-KPI estimates, season-to-season positions, and machine-learning outcomes for this inventory are reported in the Supplementary Information (Figures~S2--S3; \texttt{pef\_landscape\_2season.csv}, \texttt{ml\_empirical\_results.csv}).

Two reporting roles are distinguished. \textbf{Landscape characterisation} uses the full inventory to map where invasion-game KPIs populate the $(\kappa,\rho)$ plane, to summarise domain-level distributions, and to provide context for Figures~\ref{fig:pef_landscape}--\ref{fig:pef_ml}. These aggregates are descriptive: the inventory is Q3-heavy ($\PEFlandscapeQthreePct\%$ of KPIs), so pooled statistics primarily reflect the dominant negative-correlation regime rather than a balanced sample of the quadrant space. \textbf{Mechanistic confirmation} uses the idealised probit simulation (Tier~1) together with four curated exemplars---one per quadrant---selected at broadly comparable $\delta/\sigma_{\mathrm{A}}$ to test directional predictions of the taxonomy (\cref{sec:exemplars,tab:exemplars}).

Where domain summaries report a \textbf{landscape efficiency rate}, this is the proportion of KPI studies with $\hat\eta>1$ (pairing reduces variance relative to independent measurement). This is a descriptive summary of the $(\kappa,\rho)$ inventory, not a hypothesis-test outcome; quadrant-level aggregates appear in Supplementary Table~S1.

\subsection{Machine Learning Validation}

To assess the relationship between quadrant position and ML performance, logistic regression was fitted with five-fold cross-validation \parencite{kohavi1995}, implemented in {MATLAB} (\texttt{fitglm}, Statistics and Machine Learning Toolbox). Performance improvement ($\Delta\mathrm{ML}$) was quantified as the accuracy gain from relative features ($X_{\mathrm{A}}-X_{\mathrm{B}}$) versus absolute features ($X_{\mathrm{A}}$ only). Signal strength was estimated as $\hat{\delta}/\hat{\sigma}_{\mathrm{A}} = (\bar{X}_{\mathrm{A}}-\bar{X}_{\mathrm{B}})/s_{\mathrm{A}}$ (\cref{eq:delta_ratio}).

\subsection{Statistical Rigour}

\textbf{Multiple comparison corrections.} Where hypothesis tests were applied across a fixed battery of sports KPI checks, three methods were considered: Bonferroni correction (family-wise error rate at $\alpha=0.05$), Benjamini--Hochberg FDR control \parencite{benjamini1995}, and Holm--Bonferroni step-down procedure \parencite{holm1979}. The primary sports analysis reports per-KPI bootstrap intervals (\S\ref{sec:results}, Supplementary Information).

\textbf{Effect sizes.} Cohen's $d=(\mu_{\mathrm{relative}}-\mu_{\mathrm{absolute}})/\sigma_{\mathrm{pooled}}$ \parencite{cohen1988} and partial eta-squared $\eta_{\mathrm{p}}^{2}=\mathrm{SS}_{\mathrm{effect}}/(\mathrm{SS}_{\mathrm{effect}}+\mathrm{SS}_{\mathrm{error}})$ quantified practical significance.

\textbf{Bootstrap inference.} 10{,}000 resamples provided distribution-free confidence intervals \parencite{efron1979}.

\textbf{Power analysis.} Both post-hoc observed power and \emph{a priori} required sample sizes were calculated for small ($d=0.2$), medium ($d=0.5$), and large ($d=0.8$) effect sizes \parencite{cohen1992}.

\subsection{Quality Control}
\label{sec:qc}

All datasets underwent normality testing (Shapiro--Wilk), outlier detection (Mahalanobis distance \parencite{mahalanobis1936}), and missing data analysis (Little's MCAR test). Statistical assumptions were verified through Q--Q plots, Durbin--Watson statistics, and Levene's test where applicable. Analysis code was version-controlled with fixed random seeds.

\subsection{Sports Data: Independence Considerations}

The sports datasets present a hierarchical structure requiring careful treatment. Each team appears in multiple games, inducing correlation across observations involving the same team through latent team ability $\theta_{j}$:
\begin{equation}
  X_{j,i} = \theta_{j} + \epsilon_{j,i}.
\end{equation}
We model game-level observations using a bivariate normal approximation
\begin{equation}
  \bigl(X_{\mathrm{A}(i)}, X_{\mathrm{B}(i)}\bigr) \sim \mathcal{N}(\boldsymbol{\mu},\boldsymbol{\Sigma}),
  \quad i=1,\ldots,n,
\end{equation}
treating games as conditionally independent. Three considerations justify this approximation for PEF analysis:
\begin{enumerate}
  \item \textbf{Target of inference.} PEF focuses on within-game correlation $\rho=\Corr(X_{\mathrm{A}},X_{\mathrm{B}})$, capturing shared environmental factors and competitive dynamics rather than between-game dependencies. Team quality primarily affects mean structure rather than the correlation structure driving PEF.
  \item \textbf{Variance estimation.} The total variance $\sigma^{2}_{\mathrm{A}}=\Var(\theta_{j})+\Var(\epsilon_{j,i})$ correctly reflects the observable variance relevant to feature engineering decisions.
  \item \textbf{Robustness checks.} We validate through cluster-robust standard errors \parencite{cameron2015}, team-stratified bootstrap resampling, and sensitivity analysis comparing PEF estimates under independence versus mixed-effects modelling. These checks suggest that the independence approximation does not materially affect PEF estimates.
\end{enumerate}
This parallels established practice in sports analytics \parencite{benaim2006,owen2014} and panel data econometrics \parencite{cameron2015,wooldridge2010}.
